\newpage
\begin{center}
	{\large The Stratification of Culture Talk}
\end{center}

\begin{abstract}
\small In this paper, I explicitly characterize and establish the substantive relevance of a phenomenon that is crucial in connecting cultural inequalities to stratification outcomes:  namely, ``culture talk,'' or the use of conversational content associated with popular culture and the arts in face-to-face interaction with friends and acquaintances.  I argue that this phenomenon---central to many \emph{theories} of the connection between culture and stratification---deserves more \emph{empirical} attention.  To begin to bridge this gap, I analyze survey data from a probability sample of Americans to the first time document the socio-demographic correlates of culture talk.  The results show that the frequency of culture talk across five culture consumption domains is stratified according to the traditional markers of socio-economic status (such as levels of schooling and income) net of self-reported culture-consumption frequency, that the highly educated are more likely to engage in culture talk than the lower educated at comparably (high) levels of consumption and that consumption of in traditionally high status domains (the fine arts) is a better predictor of culture talk across all domains than consumption in less status stratified domains.  
\end{abstract}